\chapter{Built-in Factory Voices}
\label{chap:factory-voices}

Yeban's built-in poly synth has six factory presets. They share one engine: two
oscillator branches, a four-pole ladder low-pass, and a fixed ADSR; the only
difference is the nine parameter values. This chapter walks through each preset
with its parameters, typical use, and boundaries, so a preset name is not mistaken
for a separate synthesis engine.

\section{The common base of the six presets}

All six use the \texttt{synth:} prefix and are reached from the left sound library
and the bottom device rack. The same preset has slightly different labels in the
left library and the arrangement; use the preset ID and engine name as the source
of truth. \begin{center}
\begin{tikzpicture}[
  node distance=8mm and 30mm,
  core/.style={draw, rounded corners=2pt, text width=54mm, minimum height=10mm, align=center, font=\small, fill=black!4},
  preset/.style={draw, rounded corners=2pt, minimum width=24mm, minimum height=7mm, align=center, font=\scriptsize},
  arrow/.style={-{Latex[length=1.6mm]}, semithick}
]
  \node[core] (engine) {One engine: two oscillators + four-pole low-pass + fixed ADSR};
  \node[preset, above left=8mm and 20mm of engine] (h) {Hollow Lead};
  \node[preset, above=8mm of engine] (o) {Warm Organ};
  \node[preset, above right=8mm and 20mm of engine] (c) {Choir Pad};
  \node[preset, below left=8mm and 20mm of engine] (f) {FM Bell};
  \node[preset, below=8mm of engine] (g) {Glass Pluck};
  \node[preset, below right=8mm and 20mm of engine] (s) {Sub Bass};
  \foreach \p in {h,o,c,f,g,s} { \draw[arrow] (engine) -- (\p); }
\end{tikzpicture}
\end{center}

None of the six contains an adjustable ADSR, LFO, modulation matrix, or
separate envelope stages. There are only nine parameters:

\begin{itemize}
  \item Filter: \texttt{cutoff\_hz}, \texttt{resonance}, \texttt{drive};
  \item Two oscillator branches: \texttt{osc1\_wave}, \texttt{osc2\_wave},
        \texttt{osc1\_level}, \texttt{osc1\_detune\_cents},
        \texttt{osc2\_level}, \texttt{osc2\_detune\_cents}.
\end{itemize}

The wavetable index is fixed: 0 = hollow, 1 = organ, 2 = vocal, 3 = metallic,
4 = glass. A branch level of 0 bypasses that branch bit-exactly; detune is in
cents with a device range of ±4800. The device clamps cutoff between 20 Hz and
0.45 times the sample rate; resonance and drive are 0–1.

\section{Parameter tables for the six presets}

\begin{table}[htbp]
\centering
\small
\begin{tabular}{@{}lrrr@{}}
\hline
Preset & \texttt{cutoff\_hz} & \texttt{resonance} & \texttt{drive} \\
\hline
Hollow Lead & 2400 & 0.20 & 0.00 \\
Warm Organ & 1800 & 0.15 & 0.10 \\
Choir Pad & 2200 & 0.25 & 0.00 \\
FM Bell & 4200 & 0.30 & 0.15 \\
Glass Pluck & 3600 & 0.20 & 0.00 \\
Sub Bass & 420 & 0.35 & 0.25 \\
\hline
\end{tabular}
\caption{Filter parameters of the six factory presets}
\label{tab:voice-filter}
\end{table}

\begin{table}[htbp]
\centering
\small
\begin{tabular}{@{}lrrrrrr@{}}
\hline
Preset & \texttt{osc1\_wave} & \texttt{osc2\_wave} & \texttt{osc1\_level} & \texttt{osc1\_detune} & \texttt{osc2\_level} & \texttt{osc2\_detune} \\
\hline
Hollow Lead & 0 & 0 & 0.90 & $-4$ & 0.35 & $+5$ \\
Warm Organ & 1 & 1 & 0.80 & $-5$ & 0.60 & $+6$ \\
Choir Pad & 2 & 0 & 0.85 & $-8$ & 0.55 & $+8$ \\
FM Bell & 3 & 4 & 0.85 & $-3$ & 0.50 & $+11$ \\
Glass Pluck & 4 & 4 & 0.85 & $-4$ & 0.45 & $+4$ \\
Sub Bass & 0 & 1 & 0.95 & $0$ & 0.35 & $-5$ \\
\hline
\end{tabular}
\caption{Oscillator parameters of the six factory presets; detune in cents}
\label{tab:voice-osc}
\end{table}

\screenshot{factory-preset-picker}{The preset entry in the bottom device rack is open, all six synth presets are visible, and the current track is a MIDI synth track.}

\section{Hollow Lead}

Hollow Lead is the historical default: the first branch uses the hollow table at
level 0.90 and −4 cents detune; the second branch also uses hollow at level 0.35
and +5 cents; cutoff 2400 Hz, resonance 0.20, drive 0.

\textbf{Use}: soft lead lines and short melodic phrases, or one layer inside a bed.
Two branches of the same table, slightly detuned, sound thicker than one branch
without turning into a different instrument.

\textbf{Boundaries}: it is one synth voice, not a complete chord bed; filling out
harmony still needs multiple tracks or arrangement. It has no adjustable ADSR;
attack, decay, sustain, and release are fixed at 5 ms / 80 ms / 0.7 / 50 ms.

\section{Warm Organ}

Warm Organ uses the organ table on both branches: the first at level 0.80 and
−5 cents; the second at level 0.60 and +6 cents; cutoff 1800 Hz, resonance 0.15,
drive 0.10.

\textbf{Use}: organ-like sustained tones, chord pads, and passages that need a
stable pitch. Both organ branches are open and detuned for thickness; drive 0.10
adds a little harmonic content without changing overall loudness.

\textbf{Boundaries}: it has no organ drawbars, rotary speaker, or percussion key
control; ``Organ'' describes a timbre direction. Those effects need separate
effects, not a hidden control in this preset.

\section{Choir Pad}

Choir Pad uses vocal on the first branch at level 0.85 and −8 cents; hollow on the
second at level 0.55 and +8 cents; cutoff 2200 Hz, resonance 0.25, drive 0.

\textbf{Use}: vocal-coloured beds, ambient sections, and harmonic backgrounds.
The vocal table supplies the main colour while the hollow branch adds body;
the wider detune keeps the two branches from lining up exactly.

\textbf{Boundaries}: vocal is a synthesis wavetable, not a sampled choir, and there
is no formant, lyric, or phoneme control. The labels ``Vocal Lead / Choir Pad''
describe a timbre, not a singing synthesizer. For real voices, use audio-track
material.

\section{FM Bell}

FM Bell uses metallic on the first branch at level 0.85 and −3 cents; glass on the
second at level 0.50 and +11 cents; cutoff 4200 Hz, resonance 0.30, drive 0.15.

\textbf{Use}: bright, metallic, bell-like short phrases and accents. Two high
tables, a high cutoff, and a little drive produce a bright, hard transient.

\textbf{Boundaries}: the ``FM'' in the name is only preset naming; Yeban's built-in
synth is not an FM engine. It is a two-oscillator wavetable synth with a four-pole
ladder low-pass, with no carrier/modulator ratio, FM depth, or operator envelope
controls. Do not apply an FM synth parameter chart here.

\section{Glass Pluck}

Glass Pluck uses glass on both branches: the first at level 0.85 and −4 cents; the
second at level 0.45 and +4 cents; cutoff 3600 Hz, resonance 0.20, drive 0.

\textbf{Use}: plucked, glassy, short point-like sounds and bright ornaments. The
slight detune between two branches of the same table keeps the pluck from sounding
thin.

\textbf{Boundaries}: it has no separate pluck decay, string damping, or per-note
envelope control; the envelope is still the fixed ADSR. For a shorter or longer
tail, use arrangement, note duration, or later effects; there is no pluck-decay
control in this preset.

\section{Sub Bass}

Sub Bass uses hollow on the first branch at level 0.95 and 0 cents; organ on the
second at level 0.35 and −5 cents; cutoff 420 Hz, resonance 0.35, drive 0.25.

\textbf{Use}: low-end foundation, sub bass, and sections that need a heavy bottom.
The 420 Hz low-pass removes the highs and leaves the low end; the organ branch
adds a little low-mid body.

\textbf{Boundaries}: it has no separate sub oscillator, octave switch, or glide
control; the low end comes from the low cutoff and the two branch levels. For more
complex bass lines, solve it in MIDI arrangement and effects rather than expecting
one preset to cover the whole bass design.

\section{Preset selection and saving boundaries}

The six factory presets are parameter combinations of the built-in synth, not
instrument presets in the user library. A built-in synth track has no portable
instrument identity and cannot be saved as an instrument preset through ``Save
current track as preset''. To keep a whole-track setup, use a track snapshot or
``Save as Template''. Chapter 1 covered this boundary; this chapter does not repeat
it in detail.

\section{Help-topic mapping for this chapter}

\begin{table}[htbp]
\centering
\small
\begin{tabular}{@{}p{0.34\linewidth}p{0.52\linewidth}@{}}
\hline
Section & In-app help anchor \\
\hline
The common base of the six presets & \texttt{voice-common-base} \\
Parameter tables for the six presets & \texttt{voice-parameter-tables} \\
Hollow Lead & \texttt{voice-hollow-lead} \\
Warm Organ & \texttt{voice-warm-organ} \\
Choir Pad & \texttt{voice-choir-pad} \\
FM Bell & \texttt{voice-fm-bell} \\
Glass Pluck & \texttt{voice-glass-pluck} \\
Sub Bass & \texttt{voice-sub-bass} \\
Preset selection and saving boundaries & \texttt{voice-preset-boundary} \\
\hline
\end{tabular}
\caption{Help-topic mapping for Chapter 3}
\end{table}
